\subsection{IMAGE OF A COMPACT SPACE UNDER A CONTINUOUS MAPPING}

THEOREM 2. If \(f\) is a continuous mapping of a quasi-compact space \(\mathbf{X}\) into a topological space \(\mathbf{X}'\) , then the set \(f(\mathbf{X})\) is quasi-compact.

Let \(\Re\) be a covering of \(f(\mathbf{X})\) by open sets in \(\mathbf{X}'\) ; then \(\overline{f}^{1}(\Re)\) is an open covering of \(\mathbf{X}\) ( \(\S 2\) , no. 1, Theorem 1); hence there is a finite subset \(\mathfrak{S}\) of \(\Re\) such that \((\mathfrak{S}\overline{f}^{1})\) is a covering of \(\mathbf{X}\) ; but then \(\mathfrak{S}\) is a covering of \(f(\mathbf{X})\) and the theorem is proved.

COROLLARY I. Let \(f\) be a continuous mapping of a topological space \(\mathbf{X}\) into a Hausdorff space \(\mathbf{X}'\) . Then the image under \(f\) of any quasi-compact (resp. relatively quasi-compact) set in \(\mathbf{X}\) is a compact (resp. relatively compact) set in \(\mathbf{X}'\) .

COROLLARY 2. Every continuous mapping \(f\) of a quasi-compact space \(X\) into a Hausdorff space \(X'\) is closed. If also \(f\) is bijective, then \(f\) is a homeomorphism.

This follows immediately from Corollary 1 and Proposition 4 of no. 3.

In particular :

COROLLARY 3. A Hausdorff topology which is coarser than the topology of a quasi-compact space must coincide with the latter.

COROLLARY 4. Let X be a topological space and R a Hausdorff equivalence relation on X.

\begin{enumerate}
\def\labelenumi{\alph{enumi})}
\item
  If there is a quasi-compact set K in X which meets every equivalence class mod R, then X/R is compact and the canonical mapping of \(K/R_{K}\) onto X/R is a homeomorphism.
\item
  If K also meets each equivalence class in only one point, then K is a continuous section of X with respect to the relation R (§ 3, no. 5).
\end{enumerate}

Let \(f\) be the restriction to \(K\) of the canonical mapping \(X \to X / R\) . Since \(X / R\) is Hausdorff it follows from Corollary 1 that \(X / R\) is compact and from Corollary 2 that \(f\) is closed; hence the bijection \(K / R_K \to X / R\) associated with \(f\) is a homeomorphism (§ 5, no. 2, Proposition 3). This deals with a); b) follows immediately, since we now have \(K / R_K = K\) .

